\documentclass[svgnames, tikz]{standalone}
\usetikzlibrary{positioning}
\usetikzlibrary[shapes.geometric]
\usetikzlibrary {shapes.misc}
\usetikzlibrary {decorations.text}

\usetikzlibrary{calc}


\usepackage{fontspec}
\newfontfamily\Chinese{msyh.ttc}

\begin{document}

\tikzset{Arrow/.style={->, >=stealth, draw}}

\tikzset{Terminal/.style={draw=Green, line width=1pt, shape=rounded rectangle, inner sep=0.8em}}
\tikzset{Decision/.style={draw, shape=diamond, aspect=1.3, inner sep=0pt, minimum size=3em}}

\begin{tikzpicture}


\node[Decision] (ck)  at (0,0) {-c + V-};

\node[right=of ck, Terminal] (ck-end) { -ckV- };

\path[Arrow, postaction={decorate, decoration={raise=2pt, text align=center, text along path, text={Yes}}}] (ck) -- (ck-end);

\node[left=of ck,  Decision] (yi) {-Vy + ? } ;
\node[left=of yi, Terminal, align=center] (yi-end) { -V\textcolor{red}{i}? \\ \Chinese 注意拼写规则 };
\path[Arrow, postaction={decorate, decoration={raise=2pt, reverse path, text align=center, text along path, text={Yes}}}] (yi) -- (yi-end);

\coordinate (p1)  at ([yshift=1cm]$(ck)!0.5!(yi)$);
\node[above=of p1, Terminal] (start) {};
\draw (start) -- (p1);
\draw[Arrow] (p1) -| (ck);
\draw[Arrow] (p1) -| (yi);

\coordinate (p2)  at ([yshift=-1cm]$(ck)!0.5!(yi)$);

\node[below=of p2, Decision, align=center] (consecutive3) {-XX + X- \\ \Chinese 连续三个字母相同};
\node[right=of consecutive3, Decision, align=center] (consecutive3-1) {\Chinese 后缀只有一个字母};
\node[right=of consecutive3-1, Terminal, align=center] (consecutive3-1-end) {-XXeX \\ \Chinese 中间补e};
\path[Arrow, postaction={decorate, decoration={raise=2pt, text align=center, text along path, text={Yes}}}] (consecutive3) -- (consecutive3-1);
\path[Arrow, postaction={decorate, decoration={raise=2pt, text align=center, text along path, text={Yes}}}] (consecutive3-1) -- (consecutive3-1-end);

\node[below=of consecutive3-1, Terminal, align=center] (consecutive3-more) {-XX- \\ \Chinese 删除一个相同字母};
\path[Arrow, postaction={decorate, decoration={raise=2pt, text align=center, text along path, text={No}}}] (consecutive3-1) -- (consecutive3-more);

\draw (yi) |- (p2);
\draw (ck) |- (p2);
\node[anchor=south, yshift=2pt] at (p2) {No};
\draw[Arrow] (p2) -- (consecutive3);

\node[below=of consecutive3, draw, inner sep=0.8em, align=center] (drop-e) {Drop E rule 
\\ \Chinese 记住原词和变形后的词
%\\\Chinese 记住e是否被删去
};
\draw[Arrow, postaction={decorate, decoration={raise=2pt, text align=center, text along path, text={No}}}] (consecutive3) -- (drop-e);

\node[below=of drop-e, Decision, align=center] (ending-sound-v) {\Chinese 原词\\
\Chinese 结尾发音是元音};


\node[below=of ending-sound-v, Decision, align=center] (ending-sound-cc) {\Chinese 原词结尾发音\\\Chinese 是连续辅音};

\node[below=of ending-sound-cc, Decision, align=center] (ending-letter-v) {\Chinese 变形后\\
\Chinese 结尾字母是元音字母
\\(or w)};


\node[below=of ending-letter-v, Decision, align=center] (ending-letters) {
\Chinese 变形后\\
\Chinese 结尾字母是\\ h,k,q,v,w,x,y,z\\ or c};


\node[below=of ending-letters, anchor=east, align=center] (hint) {\Chinese 至此，原词结尾发音一定是辅音，
\\\Chinese 结尾辅音前一个音一定是元音
\\\Chinese 变形后的结尾字母一定是可双写字母};

\draw[Arrow] (drop-e) -- (ending-sound-v);
\draw[Arrow, postaction={decorate, decoration={raise=2pt, text align=center, text along path, text={No}}}] (ending-sound-v) -- (ending-sound-cc);
\draw[Arrow, postaction={decorate, decoration={raise=2pt, text align=center, text along path, text={No}}}] (ending-sound-cc) -- (ending-letter-v);
\draw[Arrow, postaction={decorate, decoration={raise=2pt, text align=center, text along path, text={No}}}] (ending-letter-v) -- (ending-letters);
%\draw[Arrow, postaction={decorate, decoration={raise=2pt, text align=center, text along path, text={No}}}] (ending-letters) -- (hint);

\coordinate (p3) at ([xshift=1cm]ending-letter-v.east);
\node[right=of p3, Terminal] (simple-connect) {Simple connect};

\draw[postaction={decorate, decoration={raise=2pt, text along path, text={Yes}}}] (ending-sound-v) -| (p3);
\draw[postaction={decorate, decoration={raise=2pt, text along path, text={Yes}}}] (ending-sound-cc) -| (p3);
\draw[postaction={decorate, decoration={raise=2pt, text along path, text={Yes}}}] (ending-letter-v) -| (p3);
\draw[postaction={decorate, decoration={raise=2pt, text along path, text={Yes}}}] (ending-letters) -| (p3);
\draw[Arrow] (p3) -- (simple-connect);

\node[below=2cm of ending-letters, Decision, align=center] (long-vowel) {\Chinese 原词结尾辅音\\
\Chinese 前一个音是长元音 };

\node[below=of long-vowel, Decision, align=center] (more-syllables) {\Chinese 原词音节多于一个};

\draw[Arrow, postaction={decorate, decoration={raise=2pt, text align=center, text along path, text={No}}}] (ending-letters) -- (long-vowel);
\draw[Arrow, postaction={decorate, decoration={raise=2pt, text align=center, text along path, text={No}}}] (long-vowel) -- (more-syllables);
\draw[postaction={decorate, decoration={raise=2pt, text along path, text={Yes}}}] (long-vowel) -| (p3);
\draw[postaction={decorate, decoration={raise=2pt, text along path, text={Yes}}}] (more-syllables) -| (p3);

\node[below=of more-syllables, Terminal] (doubling) {\Chinese 双写变形后的结尾字母};
\draw[Arrow, postaction={decorate, decoration={raise=2pt, text align=center, text along path, text={No}}}] (more-syllables) -- (doubling);

\end{tikzpicture}


\end{document}
